\section{Move the six-block plan into assembly}
The preceding kernel receives an already positioned eight-by-eight window.
The next routine receives complete current and reference planes, a macroblock
position and the selected motion displacement. It must derive every block origin
itself. This removes the host preparation helper from the geometry calculation
that the finished encoder will perform.

Implement \texttt{h263\_predict\_mb}. Pass an eight-byte-aligned 144-byte request
in \texttt{\$a0} and a sixteen-byte-aligned 1152-byte output area in
\texttt{\$a1}. Request fields are native-endian 64-bit values. The displacement
components are signed; all other fields are unsigned.
\begin{center}
  \tagpdfsetup{table/header-rows={1}}
  \begin{tabular}{rll}
    \textbf{Byte offsets} & \textbf{Fields in order}     & \textbf{Meaning}                      \\
    0, 8, 16              & Current Y, Cb, Cr pointers   & Complete plane origins                \\
    24, 32, 40            & Reference Y, Cb, Cr pointers & Complete plane origins                \\
    48, 56, 64            & Current Y, Cb, Cr strides    & Independent byte strides              \\
    72, 80, 88            & Reference Y, Cb, Cr strides  & Independent byte strides              \\
    96, 104               & $x$, $y$                     & Current luma macroblock origin        \\
    112, 120              & $d_x$, $d_y$                 & Selected full-luma-pixel displacement \\
    128                   & Mode                         & 0 scalar; 1 MSA block arithmetic      \\
    136                   & Reserved                     & Zero
  \end{tabular}
\end{center}
Within each three-field group, the order is Y, Cb, Cr. Y strides are 176 through
4096; chroma strides are 88 through 4096. Do not infer one plane's stride from
another's. The caller supplies all active samples of the six planes: each Y
plane is 176 by 144, and each Cb or Cr plane is 88 by 72. Byte pointers may be
unaligned. Read-only inputs may overlap. The output is disjoint from the inputs
and request. Supply a sixteen-byte-aligned stack pointer and valid private stack
space; the routine restores that incoming pointer before returning.

The current origin must be a macroblock origin: both coordinates are multiples
of sixteen, with $0\leq x\leq160$ and $0\leq y\leq128$. Each displacement is
between $-2$ and $2$, inclusive. Require the entire selected reference luma block
to remain inside the picture: $0\leq x+d_x\leq160$ and
$0\leq y+d_y\leq128$. Reject an out-of-picture choice; this routine does not
repeat the earlier search or clamp the supplied displacement.

Validate record alignments and non-wrapping extents before request reads. Then
validate all fields and all six complete plane extents before any pixel read.
For a plane of width $w$, height $h$ and stride $s$, the minimal extent is
$(h-1)s+w$. This check does not permit reading padding. Actual reads are restricted
to the current macroblock's active samples and the selected reference blocks
with their required interpolation halos.

For argument rejection, return \texttt{\$v0}=1 and \texttt{\$v1}=0, with no
pixel reads or output changes. If a nested block predictor returns nonzero,
return status two and sum zero, retaining every caller output byte. Earlier
block work may already have read pixels in that case. Development status four
remains the supplied unfinished routine's convention.

On success, return status zero and the sum of all six block SAD values.
The upper bound is $6\times16\,320=97\,920$. Store six 192-byte prediction/residual
records in the following order, regardless of your processing order:
\begin{center}
  \tagpdfsetup{table/header-rows={1}}
  \begin{tabular}{rrl}
    \textbf{Block} & \textbf{Record start} & \textbf{Contents} \\
    0              & 0                     & Y top-left        \\
    1              & 192                   & Y top-right       \\
    2              & 384                   & Y bottom-left     \\
    3              & 576                   & Y bottom-right    \\
    4              & 768                   & Cb                \\
    5              & 960                   & Cr
  \end{tabular}
\end{center}
Each record retains the previous kernel's layout: 64 predicted bytes followed
by 128 bytes of signed halfwords. Thus all six records contain 384 prediction
bytes and 768 residual bytes. These are working arrays for subsequent transforms,
not an H.263 serialization of a macroblock.

Preserve saved general registers, the stack pointer, the request and all input
bytes. Reuse your completed \texttt{h263\_predict8} or another equivalent exact
implementation. Public feedback accepts different block orders, private frame
layouts and register choices when the observable contract agrees.
